\subsection{Associated fiber bundles}

Let E and F be sets, and let G be a group acting on the right on E and on the left on F. The group G acts on the right on the product \(E \times F\) by the law \((x, y) \cdot g = (x \cdot g, g^{-1} \cdot y)\) , for \(g \in G\) and \((x, y) \in E \times F\) . The quotient set of \(E \times F\) by this operation is denoted \(E \times^{G} F\) .

When E and F are topological spaces, one equips the set \(E \times^{G} F\) with the quotient topology induced by that of \(E \times F\) . The canonical map of \(E \times F\) onto \(E \times^{G} F\) is continuous. If the group G acts continuously on E and F, it is open.

Additionally, let \(F'\) be a set on which G acts on the left, and let \(h: F \to F'\) be a map compatible with the operations of G on F and \(F'\) (A, I, p. 50). The map \(Id_{E} \times h: E \times F \to E \times F'\) is compatible with the operations of G and, by passing to quotients, defines a map denoted \(Id_{E} \times^{G} h\) from \(E \times^{G} F\) into \(E \times^{G} F'\) .

If \(h \colon F \to F'\) is a continuous map (compatible with the operations of G on F and \(F'\) ), the map \(\mathrm{Id}_{\mathrm{E}} \times^{\mathrm{G}} h\) is continuous (TG, I, p. 21, prop. 6).

Examples. --- 1) Let F be a topological space and let G be a topological group acting continuously on the left on F. If one endows the topological space G with the action of G by right translations, the space G × \(^{G}\) F is canonically identified with F as follows. The continuous maps \(\varphi\) : F → G × F and \(\psi\) : G × F → F defined by \(\varphi(f) = (e, f)\) (where e denotes the identity element of G) and \(\psi(g, f) = g \cdot f\) induce continuous maps \(\overline{\varphi}\) : F → G × \(^{G}\) F and \(\overline{\psi}\) : G × \(^{G}\) F → F which are inverses of each other.

\begin{enumerate}
\def\labelenumi{\arabic{enumi})}
\setcounter{enumi}{1}
\item
  Example 1 generalizes as follows. Let B and F be topological spaces and let G be a topological group acting continuously on the left on F. By passing to quotients, the map from B × F to (B × G) × F given by (b, f) ↦ ((b, e), f) and the map from (B × G) × F to B × F given by ((b, g), f) ↦ (b, gf) define B-isomorphisms B × F → (B × G) × \(^{G}\) F and (B × G) × \(^{G}\) F → B × F which are inverses of each other.
\item
  Similarly, let E be a topological space and let G be a topological group acting continuously on the right on E. If the topological space G is endowed with the action of G by left translations, the space \(E \times^{G} G\) is identified with E.
\end{enumerate}

Let E and F be topological spaces, and let G be a topological group acting continuously on the right on E and on the left on F. Let B be a topological space and \(p\colon E \to B\) a continuous map such that \(p(x \cdot g) = p(x)\) for \(x \in E\) and \(g \in G\) . The map \(p \circ pr_{1}\colon E \times F \to B\) defines, by passing to the quotient, a continuous map \(p^{F}\colon E \times^{G} F \to B\) and the canonical map \(\pi\colon E \times F \to E \times^{G} F\) is a B-morphism.

Let B' be a topological space and let \(h\colon B'\to B\) be a continuous map. The group G acts continuously on the right on \(B'\times_{B}E\) by the operation law \(((b',x),g)\mapsto(b',x\cdot g)\). By composing the canonical isomorphism \(((b',x),y)\mapsto(b',(x,y))\) of \((\mathrm{B}'\times_{\mathrm{B}}\mathrm{E})\times\mathrm{F}\) onto \(\mathrm{B}'\times_{\mathrm{B}}(\mathrm{E}\times\mathrm{F})\) and the map \(\operatorname{Id}_{\mathrm{B}'}\times\pi\) from \(\mathrm{B}'\times_{\mathrm{B}}(\mathrm{E}\times\mathrm{F})\) into \(\mathrm{B}'\times_{\mathrm{B}}(\mathrm{E}\times^{\mathrm{G}}\mathrm{F})\), we obtain a continuous map \(\lambda_{0}\colon(\mathrm{B}'\times_{\mathrm{B}}\mathrm{E})\times\mathrm{F}\to\mathrm{B}'\times_{\mathrm{B}}(\mathrm{E}\times^{\mathrm{G}}\mathrm{F})\). By passing to the quotient, it defines a continuous map

\[
\lambda \colon (\mathrm{B} ^ {\prime} \times_ {\mathrm{B}} \mathrm{E}) \times^ {\mathrm{G}} \mathrm{F} \rightarrow \mathrm{B} ^ {\prime} \times_ {\mathrm{B}} (\mathrm{E} \times^ {\mathrm{G}} \mathrm{F}).
\]

Lemma. --- The map λ is a homeomorphism.

The map \(\lambda_0\) is surjective, and two elements of \((\mathrm{B}' \times_{\mathrm{B}} \mathrm{E}) \times \mathrm{F}\) have the same image under \(\lambda_0\) if and only if they belong to the same orbit for the action of G on \((\mathrm{B}' \times_{\mathrm{B}} \mathrm{E}) \times \mathrm{F}\). It follows that the map \(\lambda\) is bijective. Since the map \(\pi\) is open, so is the map \(\mathrm{Id}_{\mathrm{B}'} \times_{\mathrm{B}} \pi\) (I, p. 17, prop. 8), hence so are \(\lambda_0\) and \(\lambda\) (TG, I, p. 32, prop. 3), which proves that \(\lambda\) is a homeomorphism.

PROPOSITION 7. --- Let B be a topological space, let G be a topological group, let (E,p) be a principal fiber bundle over B with group G, and let F be a topological space on which G acts continuously on the left.

\begin{enumerate}
\def\labelenumi{\alph{enumi})}
\item
  The topological space \(E \times^{G} F\) together with the continuous map \(p^{F}: E \times^{G} F \to B\) deduced from the map \(p \circ pr_{1}: E \times F \to B\) by passing to the quotient is a locally trivial fiber bundle of fiber type F; it is trivializable if the B-bundle E is trivializable.
\item
  Let \(\pi\colon\mathrm{E}\times\mathrm{F}\to\mathrm{E}\times^{\mathrm{G}}\mathrm{F}\) be the canonical surjection. The map \(\mu\colon\mathrm{E}\times\mathrm{F}\to\mathrm{E}\times_{\mathrm{B}}(\mathrm{E}\times^{\mathrm{G}}\mathrm{F})\) which associates \((x,\pi(x,f))\) with \((x,f)\) is a homeomorphism whose inverse map is a trivialization of the locally trivial fiber bundle over E \((\mathrm{E}\times_{\mathrm{B}}(\mathrm{E}\times^{\mathrm{G}}\mathrm{F}),\mathrm{pr}_{1})\).
\end{enumerate}

Suppose first that E is the trivial principal fiber bundle \(B \times G\) . By Example 2, the space \(E \times^{G} F\) is then identified with \(B \times F\) and the map \(p^{F}\) with the first projection \(pr_{1}: B \times F \to B\) , which proves that \((\mathrm{E} \times^{\mathrm{G}}\mathrm{F}, p^{\mathrm{F}})\) is a trivializable fiber bundle over B in this case. In the general case, every point of B has a neighborhood U such that the map \(p_{\mathrm{U}}: p^{-1}(\mathrm{U}) \to \mathrm{U}\) makes \(p^{-1}(\mathrm{U})\) a trivializable principal fiber bundle over U. By the preceding, \((p^{-1}(\mathrm{U}) \times^{\mathrm{G}}\mathrm{F} \to \mathrm{U}, (p_{\mathrm{U}})^{\mathrm{F}})\) is a trivializable fiber bundle over U. By the lemma above, applied to the case where \(B' = U\) and \(h: U \to B\) is the canonical injection, the same holds for the U-space \(((p^{\mathrm{F}})^{-1}(\mathrm{U}), (p^{\mathrm{F}})_{\mathrm{U}})\) deduced from \((\mathrm{E} \times^{\mathrm{G}}\mathrm{F}, p^{\mathrm{F}})\) by restriction over U. This proves that \((\mathrm{E} \times^{\mathrm{G}}\mathrm{F}, p^{\mathrm{F}})\) is a locally trivial fiber bundle over B and concludes the proof of assertion a).

The map \(\theta\colon\mathrm{E}\times\mathrm{G}\to\mathrm{E}\times_{\mathrm{B}}\mathrm{E}\) defined by \(\theta(x,g)=(x,x\cdot g)\) is a homeomorphism (I, p. 94, prop. 3) compatible with the operations of G on E \(\times\) G and on E \(\times_{\mathrm{B}}\) E given by \(((x,g),g')\mapsto(x,gg')\) and \(((x,y),g')\mapsto(x,yg')\) respectively. By passing to quotients, \(\theta\) therefore induces a homeomorphism \(\theta'\) from (E \(\times\) G) \(\times^{\mathrm{G}}\) F onto (E \(\times_{\mathrm{B}}\) E) \(\times^{\mathrm{G}}\) F. The map \(\mu\) is the composition of the homeomorphism E \(\times\) F \(\rightarrow\) (E \(\times\) G) \(\times^{\mathrm{G}}\) F (example 2), of \(\theta'\) and of the canonical homeomorphism (E \(\times_{\mathrm{B}}\) E) \(\times^{\mathrm{G}}\) F \(\rightarrow\) E \(\times_{\mathrm{B}}\) (E \(\times^{\mathrm{G}}\) F) (I, p. 105, lemma), whence \(b\) ).

Under the hypotheses of Proposition 7, the locally trivial fiber bundle (E × \(^{G}\) F, p \(^{F}\) ) is called the locally trivial fiber bundle of fiber type F associated with the principal fiber bundle (E, p). All fibers of the B-space E × \(^{G}\) F are homeomorphic to the space F. In particular, if the space F is discrete, the map p \(^{F}\) is a covering.

Examples. --- 4) Let B be a topological space, G a topological group, and let (E, p) be a principal G-bundle over B. Let H be a subgroup of G.

Denote by \(\varphi\colon\mathrm{E}\to\mathrm{E}\times(\mathrm{G}/\mathrm{H})\) and \(\psi\colon\mathrm{E}\times(\mathrm{G}/\mathrm{H})\to\mathrm{E}/\mathrm{H}\) the maps defined by \(\varphi(x)=(x,\mathrm{H})\) and \(\psi(x,g\mathrm{H})=(x\cdot g)\mathrm{H}\) . They are compatible with the projections to B and with the operations of G, and they define, by passing to quotients, morphisms of B-spaces \(\overline{\varphi}\colon\mathrm{E}/\mathrm{H}\to\mathrm{E}\times^{\mathrm{G}}(\mathrm{G}/\mathrm{H})\) and \(\overline{\psi}\colon\mathrm{E}\times^{\mathrm{G}}(\mathrm{G}/\mathrm{H})\to\mathrm{E}/\mathrm{H}\) , inverses of each other. We say that \(\overline{\varphi}\) is the canonical homeomorphism from E/H onto \(\mathrm{E}\times^{\mathrm{G}}(\mathrm{G}/\mathrm{H})\) . In particular, the topological space E/H equipped with the continuous map \(p_{\mathrm{H}}\colon\mathrm{E}/\mathrm{H}\to\mathrm{B}\) is a locally trivial fiber bundle over B with fiber type G/H.

If moreover H is a normal subgroup of G, the action of G endows, by passing to quotients, the B-space E/H with the structure of a principal fiber bundle with group G/H.

In particular, if E is a principal covering of B with group G, then E/H is a covering of B; if H is normal in G, it is a principal covering with group G/H.

\begin{enumerate}
\def\labelenumi{\arabic{enumi})}
\setcounter{enumi}{4}
\item
  Let B be a topological space, let G be a topological group, and let E be a principal G-bundle over B. Let F be a topological homogeneous space with respect to G (TG, III, p. 12). Let y be a point of F, and let \(G_{y}\) be its stabilizer. The map \(\varphi_{y}\colon x\mapsto(x,y)\) from E to \(E\times F\) defines, by passing to the quotient, a homeomorphism \(\overline{\varphi}_{y}\) of \(E/G_{y}\) onto \(E\times^{G}F\) . When the group G is abelian, the subgroup \(G_{y}\) does not depend on the point y, but the homeomorphism \(\overline{\varphi}_{y}\) in general depends on it.
\item
  Let B be a topological space, let G be a topological group, let \((E,p)\) be a principal G-bundle over B. Let H be a topological group and let \(f:G\to H\) be a morphism of topological groups; endow H with the left action of G given by \(g\cdot h=f(g)h\) .
\end{enumerate}

Let \(q \colon \mathrm{E} \times \mathrm{H} \to \mathrm{E} \times^{\mathrm{G}} \mathrm{H}\) be the canonical map; it is open, hence universally strict (I, p. 20, corollary 11). The map \(m \colon (x, h, h') \mapsto q(x, hh')\) from \(\mathrm{E} \times \mathrm{H} \times \mathrm{H}\) into \(\mathrm{E} \times^{\mathrm{G}} \mathrm{H}\) is continuous. Let \((x, h) \in \mathrm{E} \times \mathrm{H}\) , \(h' \in \mathrm{H}\) and \(g \in \mathrm{G}\) ; we have \(m(xg, f(g)^{-1}h, h') = q(xg, f(g)^{-1}hh') = q(x, hh') = m(x, h, h')\) . Consequently, there exists a unique map

\[
m ^ {\prime} \colon (\mathrm{E} \times^ {\mathrm{G}} \mathrm{H}) \times \mathrm{H} \rightarrow \mathrm{E} \times^ {\mathrm{G}} \mathrm{H}
\]

such that \(m'(q(x,h),h') = q(x,hh')\) for all \(x \in E\) and all \(h, h' \in H\) . Since q is universally strict, the map \(m'\) is continuous. This is a right action of the topological group H on the B-space \(E \times^{G} H\) .

Let U be an open subset of B such that \(E_{U}\) is isomorphic to the principal fiber bundle over U \(U \times G\) . Equipped with the operation of H, the U-space \((\mathrm{E} \times^{\mathrm{G}}\mathrm{H})_{\mathrm{U}}\) is identified with the principal fiber bundle \(U \times H\) . This proves that \(E \times^{G}H\) is a principal H-bundle over B.

DEFINITION 4. --- Let B be a topological space and let G be a topological group. A locally trivial fiber bundle X over B is said to be associated with a principal fiber bundle E over B with group G if there exists a topological space F on which the group G acts continuously on the left and a B-isomorphism from E × \(^{G}\) F onto X.

Let E be a principal G-bundle over B, and let X be a locally trivial fiber bundle over B associated with E. If the principal bundle E is trivializable, then X is trivializable (prop. 7). If B' is a topological space and \(h \colon \mathrm{B}' \to \mathrm{B}\) is a continuous map, the locally trivial fiber bundle \(\mathrm{B}' \times_{\mathrm{B}} \mathrm{X}\) deduced from X by base change is associated to the B'-principal fiber bundle \(\mathrm{B}' \times_{\mathrm{B}} \mathrm{E}\) (I, p. 105, lemma).

